\documentclass[12pt,a4paper]{article}

\input{/home/tabaka/Documents/GitHub/Defs/defs.tex}

%\pagestyle{empty} 			%usuwa nr strony

\setcounter{zadaniaRozdzial}{2}


\begin{document}
	
	\begin{center}
		\LARGE 2.2 Wzory Viete'a
	\end{center}
	
	{\large\Wzor{$$x_1\cdot x_2=-\frac{b}{a} \hspace{3cm} x_1+x_2=\frac{c}{a}$$}}
	
	\vspace{1cm}
	
	\ZbiorZadanieTwoXThree{Wyznaczyć za pomocą wzorów Viete'a wyrażenia dla funkcji kwadratowej $f(x)=zx^2+bx+c$}{$\frac{1}{x_1}+\frac{1}{x_2}$}{$x_1^2+x_2^2$}{$\frac{1}{x_1^2}+\frac{1}{x_2^2}$}{$x_1^3+x_2^3$}{$(x_1-x_2)^2$}{$\frac{x_1}{x_2}+\frac{x_2}{x_1}$}
	
	\ZbiorZadanie{Wyprowadzić wzory Viete'a.}
	
	\ZbiorZadanie{Wykzać, że $$x_1^4+x_2^4=(x_1+x_2)^4-4x_1x_2(x_1+x_2)^2+2(x_1x_2)^2$$}
	
	\ZbiorZadanie{Dla funkcji kwadratowej $f(x)=x^2+20x-19$, nie obliczająć pierwiastków tej funkcji kwadratowej, wyznaczyć
	\begin{enumerate}[a)]
		\item $x_1^2+x_2^2+2(x_1x_2+x_1+x_2)$
		\item $\frac{1}{x_1+1}+\frac{1}{x_2+1}$
		\item $\frac{x_1-9}{x_2+3}+\frac{x_2-9}{x_1+3}$
	\end{enumerate}}

	\ZbiorZadanieTwoXTwo{Wyznaczyć warunki dla których funkcja kwadratowa $f(x)=ax^2+bx+c$ ma dwa pierwiastki rzeczywiste, które}{są różnego samego znaku}{są oba dodatnie}{są oba ujemne}{różnią się conajwyżej o 6}
	
	\newpage
	
	\begin{center}
		\LARGE Zadania uzupełniające
	\end{center}
	\vspace{1.5cm}
	
	\ZbiorZadanieTwoXTwo{Przekształć poniższe wyrażenia tak, aby skorzystać ze wzorów Viete'a, a następnie je zastosować dla funkcji: $f(x)=\sqrt{3}x^2-(2\sqrt{3}+1)x+\sqrt{2}$}{$\frac{1}{x_1}+\frac{1}{x_2}$}{$(x_1-x_2)^2$}{$x_1^3x_2+2x_1^2x_2^2+x_1x_2^3$}{$\frac{4x_1+4x_2+1}{x_1}+\frac{4x_1+4x_2+1}{x_2}$}
	
	\ZbiorZadanieTwoXThree{Wiedząć, że podana funkcja kwadratowa ma dwa miejsca zerowe, wyznaczyć te pierwiastki korzystając ze wzorow Viete'a}
	{$f(x)=x^2+5x+6$}{$f(x)=x^2+7x+15$}
	{$f(x)=x^2-3x+4$}{$f(x)=x^2-7x+12$}
	{$f(x)=x^2-6x-16$}{$f(x)=-x^2+7x+18$}
	
	\ZbiorZadanieTwoXTwo{Dana jest funkcja kwadratowa $f(x)=\sqrt{2}x^2-\sqrt{6}x-\sqrt{3}$. Sprawdzić, ćzy posiada pierwiastki, a następnie korzystając ze wzorów Viete'a wyanczyć ich}{sumę kwadratów}{kwadrat sumy}{sume odwrotności}{sumę kwadratóœ odwrotności}
		
	
	
\end{document}